% !TeX spellcheck = ru_RU
\section{Обзор}
\label{sec:relatedworks}

\subsection{Обзор существующих решений}

Анализ связей в художественных текстах может выполняться вручную (например, вики-проекты с перекрёстными ссылками) или с помощью автоматических инструментов. Среди автоматических подходов выделяются:

\begin{itemize}
    \item \textbf{Общие парсеры вики} (например, \texttt{wikiextractor}) – позволяют извлечь текст, но не предоставляют встроенных механизмов для выделения сущностей и не поддерживают простую смену источника.
    \item \textbf{Системы на основе регулярных выражений} – требуют ручного составления словарей имён, что трудоёмко для обширных вселенных.
    \item \textbf{NER-модели} (spaCy, Natasha, DeepPavlov) – обучены на новостных текстах и могут давать ошибки на фандомных именах, однако для русскоязычных вики с типовыми именами модель \texttt{ru\_core\_news\_sm} демонстрирует приемлемое качество.
    \item \textbf{Графовые базы данных} (Neo4j) – удобны для хранения связей, но требуют отдельного сервера и более сложной настройки. В данной работе использована реляционная PostgreSQL.
    \item \textbf{Библиотеки для визуализации графов} (NetworkX, matplotlib, PyVis) – позволяют строить и анализировать графы совместных упоминаний.
\end{itemize}

Главное преимущество разработанной системы — возможность легко сменить фандом. В отличие от многих готовых парсеров, которые жёстко привязаны к устройству конкретного сайта, наша система позволяет взять уже написанный код для извлечения сущностей, работы с базой данных и анализа и просто подключить к нему новый парсер для другой вселенной. Всё, что нужно сделать, — это создать новый класс с указанием страниц нового фандома, а вся остальная часть проекта останется без изменений.

\subsection{Обзор используемых технологий}

При разработке системы были рассмотрены следующие технологии: Python (библиотеки Requests, BeautifulSoup для парсинга; spaCy и Natasha для NER; pandas для работы с CSV; NetworkX, matplotlib для графов), PostgreSQL (хранение данных), DataGrip (импорт). Сравнение представлено в таблице \ref{tab:comparison} в приложении А.

Итоговый выбор: Requests + BeautifulSoup – достаточная функциональность; spaCy – зарекомендовавшая себя библиотека, поддерживает нужные типы PERSON/LOC (в русской модели \texttt{ru\_core\_news\_sm} используется \texttt{PER} для персонажей); PostgreSQL – работа с DataGrip; NetworkX + matplotlib – для построения графа совместных упоминаний.

\subsection{Выводы по обзору}

Существующие системы либо не обеспечивают автоматического выделения сущностей, либо привязаны к одному источнику данных. Предлагаемое решение закрывает этот пробел за счёт модульной архитектуры и использования современных NLP-инструментов, а также добавляет визуализацию связей в виде графа.